\chapter{Threat Model}

This chapter defines the threat model used to assess oXiDe SE. It identifies
the assets to protect, the capabilities of potential adversaries, the trust
assumptions, and the attack surfaces. It then states the security objectives
pursued by the architecture and maps them to the components described in
Chapter~1. These objectives describe the intended effect of the design; they
are not presented as formally established guarantees.

%--------------------------------------------------
\section{Attacker Panorama and Attack Surfaces}

The attack surfaces considered in this chapter depend on where an adversary is
positioned in the lifecycle introduced in Chapter~1. Different actors do not
see the same interfaces, do not control the same privileges, and do not pursue
the same objectives. The present threat model therefore distinguishes three
main attack surfaces.

\begin{itemize}
\item \textbf{End-user side adversaries} interact with the secure element from
  outside. Their typical objective is to steal secrets manipulated by service
  providers or to undermine the integrity of those secrets and associated
  services. From that position, the relevant attack surfaces are the external
  APDU interface and, more broadly, the exposed hardware interface.
\item \textbf{Service-provider side adversaries} already execute code or deploy
  code in the secure element and try to violate the isolation expected by other
  service providers. Their typical objective is to obtain secrets belonging to
  competing services or to interfere with their execution. From that position,
  the relevant attack surfaces are the kernel access interface and, where the
  architecture allows it, any explicit service-sharing or inter-service
  interface exposed to peer applications.
\item \textbf{Hardware-origin adversaries} control or compromise an element
  below the software stack, such as the silicon, boot substrate, or physical
  board. Such an adversary could invalidate the execution and key-protection
  assumptions on which every higher layer depends. This class includes a
  malicious foundry, but is not modelled as an ordinary caller of the hardware
  abstraction layer. Hardware and silicon compromise are declared out of scope
  for the present software-architecture threat model.
\end{itemize}

\begin{figure}[ht]
  \centering
  \begin{tikzpicture}[
    stack/.style={
      draw,
      rounded corners=2pt,
      minimum width=11.5cm,
      minimum height=1.1cm,
      align=center,
      font=\small
    },
    sp/.style={
      draw,
      rounded corners=2pt,
      minimum width=3.2cm,
      minimum height=1cm,
      align=center,
      font=\small
    },
    label/.style={font=\small\itshape},
    >=Latex
  ]
    \node[stack, fill=white] (user) at (0,4.8) {\textbf{End User}};

    \node[sp, fill=gray!10] (sp1) at (-4.1,3.2) {\textbf{Service}\\\textbf{Provider A}};
    \node[sp, fill=gray!10] (sp2) at (0,3.2) {\textbf{Service}\\\textbf{Provider B}};
    \node[sp, fill=gray!10] (sp3) at (4.1,3.2) {\textbf{Service}\\\textbf{Provider C}};

    \draw[dashed] (-2.05,2.65) -- (-2.05,3.75);
    \draw[dashed] (2.05,2.65) -- (2.05,3.75);

    \draw[thick] (-6.0,1.6) -- (6.0,1.6);
    \node[label, anchor=west] at (-5.9,1.85) {unsafe world};
    \node[label, anchor=west] at (-5.9,1.25) {root of security};

    \node[stack, fill=gray!25] (issuer) at (0,0.0) {\textbf{Issuers}};
    \node[stack, fill=gray!70, text=white] (provider) at (0, -1.6) {\textbf{Providers}};
    \node[stack, fill=black, text=white] (foundry) at (0,-3.2) {\textbf{Foundries}};

    \draw[->, thick] (user.south) -- (0,3.75);
    \draw[->, thick] (sp2.south) -- (issuer.north);
    \draw[->, thick] (issuer.south) -- (provider.north);
    \draw[->, thick] (provider.south) -- (foundry.north);
  \end{tikzpicture}
  \caption{Trust stack around a secure element, from foundries up to end-user-facing services.}
  \label{fig:trust-stack-actors}
\end{figure}

This decomposition clarifies the reading of the rest of the chapter. The
external communication path is studied primarily because it is the current
entry point of the prototype and therefore the first concrete attack surface.
The kernel-facing and inter-service surfaces are equally important for the
target architecture, because the long-term goal is to host mutually isolated
Rustlets or applications that may originate from distinct and potentially
hostile service providers.

Not every actor category defined in Chapter~1 appears as an attacker category
here. Providers and issuers remain relevant for trust assumptions and
administrative authority, but the present chapter focuses on the interfaces
through which end users, service providers and hardware-origin adversaries can
actually attack the platform.

Figure~\ref{fig:trust-stack-actors} summarizes the actors involved in that trust
stack, but the position of an organization in the lifecycle does not by itself
make all of its software trusted. End users and service providers consume the
security properties sought by the platform. Issuers, image integrators,
provisioning authorities, toolchain operators, and hardware providers form the
pre-issuance chain of trust to the extent that they can alter or approve the
initial image and its trust anchors. The service-provider layer is intentionally
shown as multiple parallel boxes because the target architecture must support
coexistence of mutually isolated service-provider payloads within one secure
element, whereas the end-user layer is represented by a single box because one
deployed secure element is operated in one end-user context at a time.

An ordinary Rustlet remains outside the runtime TCB whether it is loaded before
or after issuance. Its developer does not become a trusted authority merely
because that Rustlet is pre-deployed. The authority that selects the Rustlet,
approves its inclusion, and produces the issued image is, however, part of the
pre-issuance chain of trust. Privileged kernel-resident modules and Security
Domain authorization logic have the more specific TCB status defined in
Chapter~1.

%--------------------------------------------------
\section{Scope and Assets}

This section defines the assets that must be protected by the platform. These
assets correspond to both data and control structures that are critical to the
security and integrity of the system.

\begin{itemize}
\item Integrity of the application lifecycle
\item Confidentiality and integrity of cryptographic keys
\item Isolation between applications
\item Integrity of the internal registry
\item Control over administrative channels
\end{itemize}

Across the three attack surfaces introduced above, these assets are exposed in
different ways: APDU-facing adversaries target command authorization and data
confidentiality, service-provider-facing adversaries target isolation and
kernel mediation, and hardware-origin adversaries would ultimately target the
root secrets on which all higher layers depend.

The integrity of the application lifecycle is the objective that operations
such as LOAD, INSTALL and DELETE must not be performed in an unauthorized or
inconsistent manner. Attacks targeting lifecycle management
have been demonstrated in smart card environments where logical flaws allow
inconsistent states or unauthorized installation flows
\cite{mostowski2008malicious, barbu2010javacard}.

The confidentiality and integrity objective for cryptographic keys requires
that sensitive material used for secure channels or internal services be
accessible and modifiable only through authorized interfaces. Numerous works
on Trusted Execution Environments and secure elements
highlight practical key extraction or misuse vulnerabilities when isolation or
API design is flawed \cite{shakevsky2022keymaster, busch2020trustedcore}.

Isolation between applications is intended to prevent one application from
interfering with or observing the execution of another. Violations of such
isolation have been widely studied in both Java Card and TEE contexts
\cite{mostowski2008malicious, cerdeira2022rezone}.

The integrity objective for the internal registry requires identifiers, states,
and privileges associated with applications and domains to remain consistent
and modifiable only through kernel-mediated transitions. Corruption or
manipulation of such metadata can lead to privilege
escalation or inconsistent system states, as observed in several TEE
vulnerability analyses \cite{munoz2023tee}.

Control over administrative channels is intended to ensure that only an entity
holding the required credentials and authority can cause management commands to
take effect. Replay, impersonation, message modification, and relay are distinct
communication threats: authentication and freshness can address the first
three, whereas a transparent real-time relay requires additional proximity or
channel-binding mechanisms that are not currently provided by oXiDe SE
\cite{roland2012relay, francis2011relay}.

General service availability is not treated as a separate asset in this model.
APDU processing is sequential, command buffers are bounded, and Rustlets use
statically constrained volatile resources, so an application cannot accumulate
unbounded APDUs or RAM allocations. Persistent storage has a narrower progress
objective: journal updates should preserve a previous valid state, and a
Rustlet whose serialized state remains within its assigned capacity should
remain rewritable. A genuine increase beyond that capacity may legitimately
return an out-of-memory error. Chapter~6 describes the persistence mechanism
against which this objective is assessed.

%--------------------------------------------------
\section{Adversary Model}

This section describes the capabilities of potential adversaries interacting
with the platform.

\begin{itemize}
\item External adversary with communication access
\item Malicious internal application
\item Compromised administrator
\end{itemize}

These adversary classes map directly onto the lifecycle-based panorama. The
external adversary corresponds to the end-user side attack surface. The
malicious internal application corresponds to the service-provider side attack
surface once independently deployed services are considered. The compromised
administrator overlaps with issuer- or provider-level operational authority and
therefore informs the privilege and trust boundaries around lifecycle
management.

The external adversary with communication access is able to send arbitrary APDU
commands to the platform and observe responses. This includes the ability to
inject malformed or malicious command sequences, attempt replay of previously
observed exchanges, or interact with the system through a man-in-the-middle
position on the communication channel. Such capabilities are illustrated by
relay and communication-layer attacks against secure elements and contactless
systems \cite{roland2012relay, francis2011relay, zhao2021sim}.

The malicious internal application corresponds to an application that has been
successfully installed but behaves in a hostile manner. Such an application may
attempt to escalate its privileges, access resources belonging to other
applications, or violate secure-storage boundaries, in particular by targeting
key management mechanisms or sensitive persistent data. These behaviors are
well documented in Java Card and TEE environments \cite{mostowski2008malicious,
barbu2010javacard, lanet2012handson, bouffard2012nightmare, bouffard2015cft,
busch2020trustedcore}.

The compromised administrator represents an entity that legitimately holds
administrative privileges but uses them in a malicious or unintended way. Its
capabilities depend on which authority is compromised. A supplementary
Security Domain may misuse the powers delegated over its own administrative
subtree; the kernel is intended to prevent it from escaping that subtree or
bypassing platform-wide invariants. Compromise of the issuer, image-production
authority, or configured root authority has a substantially larger impact and
may invalidate the initial policy itself. The model therefore does not treat
``administrator'' as one uniform privilege level. Similar threat patterns are
studied in TEE privilege models and their misuse
\cite{cerdeira2022rezone, munoz2023tee}.

%--------------------------------------------------
\section{Trust Assumptions}

This section states the assumptions required for the architecture to pursue its
security objectives.

\begin{itemize}
\item Trusted root of execution
\item Correct implementation of cryptographic primitives
\item Secure key storage
\item Integrity of the runtime environment
\end{itemize}

The trusted root of execution assumes that the initial boot process and trusted
code base cannot be altered by an adversary. The runtime TCB is the one defined
in Chapter~1: the target execution substrate, kernel, privileged kernel-resident
modules, relevant Security Domain authorization logic, trust anchors, and
security configuration. This assumption is common in secure element and TEE
architectures.

Before issuance, the compiler, image-construction tools, provisioning process,
image-approval authorities, and hardware supply chain must also behave as
expected. They are dependencies of the initial trusted image rather than
resident components of its runtime TCB. Ordinary Rustlets and their developers
remain outside that TCB, including when an integration authority chooses to
pre-deploy such a Rustlet.

The correct implementation of cryptographic primitives assumes that algorithms
such as AES or MAC computations are implemented without flaws and behave as
expected. However, numerous works have shown that even correct algorithms can
leak information through side channels \cite{kocher1996timing, kocher1999dpa}.

Secure key storage assumes that cryptographic material cannot be extracted or
modified outside of authorized interfaces. Violations of this assumption have
been demonstrated in practice in TEE implementations \cite{bouffard2014reverse,
shakevsky2022keymaster}.

The integrity of the runtime environment assumes that the execution context
enforces isolation and prevents arbitrary code execution outside controlled
mechanisms.  Attacks breaking such assumptions have been demonstrated through
both fault-assisted techniques, which enable control-flow deviation or type
confusion in smart card environments \cite{dubreuil2012fault,
razafindralambo2012virus}, and logical vulnerabilities in Trusted Execution
Environments, allowing privilege escalation or interface abuse
\cite{busch2024globalconfusion}.

%--------------------------------------------------
\section{Attack Surfaces}

This section identifies the main interfaces through which an adversary may interact with the system.

\begin{itemize}
\item APDU communication interface
\item Secure channel protocol
\item Administrative command set
\item Inter-application invocation mechanisms
\item Internal registry and persistence layer
\end{itemize}

These interfaces refine the three attack surfaces of Section~2.1. The APDU
communication interface and the secure channel protocol belong to the end-user
side entry points. The administrative command set, inter-application
invocation mechanisms, and registry-facing operations become especially
relevant on the service-provider side, where a hostile deployed component may
try to cross the kernel boundary or abuse an explicit sharing interface.

On the RP2350, the current Pico~2 port starts oXiDe SE in the processor's Secure
state, but it does not yet configure a complete Secure/Non-secure TrustZone
partition or expose the final virtual-Secure-Element interface. Completing that
integration is planned for the beta~1 development cycle. In that future
configuration, the Non-secure world and its operating system must be treated as
potentially hostile, and every Secure-world entry point becomes an additional
external attack surface.

The APDU communication interface is the primary entry point and therefore the
main attack surface, as all external commands are received through this
interface. Communication-level attacks such as relay or replay attacks directly
target this interface \cite{roland2012relay}.

The secure channel protocol represents a critical surface where errors in
authentication, key derivation or message protection could compromise the
entire system. Weaknesses in secure channel implementations have historically
led to severe compromises.

The administrative command set includes operations such as INSTALL, LOAD or
DELETE, which directly impact the system state and must be strictly controlled.
Logical attacks exploiting these commands have been demonstrated in smart card
environments \cite{barbu2010javacard, lanet2012handson}.

Where explicit inter-application invocation mechanisms are provided, they allow
applications to interact and may be exploited if isolation or access control is
insufficient. They are therefore a target-architecture surface, not a claim
that every current image exposes such an interface. Such weaknesses are well
documented in both Java Card and TEE contexts
\cite{mostowski2008malicious, cerdeira2022rezone}.

The internal registry and persistence layer store critical metadata and
represent a target for tampering or corruption attacks, which can lead to
persistent privilege escalation \cite{munoz2023tee}.

%--------------------------------------------------
\section{Security Goals}

This section defines the security properties that the architecture is designed
to enforce. Their implementation and validation evidence are examined in later
chapters; no formal proof is implied here.

\begin{itemize}
\item Authenticity
\item Integrity
\item Confidentiality
\item Isolation
\item Non-bypassability
\end{itemize}

Authenticity aims to ensure that only entities holding the required credentials
can issue protected administrative or operational commands. It mitigates
impersonation, but does not by itself establish physical proximity or prevent a
transparent real-time relay \cite{roland2012relay}.

Integrity aims to make modifications of system state and application code
possible only through authorized transitions. It addresses logical
modification attempts within the stated model; physical fault injection remains
out of scope \cite{boneh1997faults}.

Confidentiality aims to protect sensitive data, in particular cryptographic
keys and secure-storage contents, against unauthorized logical access.
Side-channel attacks remain out of scope \cite{kocher1999dpa}.

Isolation aims to make applications operate independently and to contain a
malicious or defective Rustlet within its mediated resources, a key requirement
highlighted in multiple Java Card and TEE attacks
\cite{mostowski2008malicious}.

Non-bypassability is the objective that privileged effects remain mediated by
the kernel and cannot circumvent the security model, a property often violated
in real-world TEE vulnerabilities \cite{cerdeira2022rezone}.

%--------------------------------------------------
\section{Mitigations and Architectural Mapping}

This section relates the identified threats to the architectural components
previously described.

\begin{itemize}
\item Secure channel protocol
\item Privilege model
\item Security Domains
\item Registry consistency mechanisms
\item Runtime isolation
\end{itemize}

The mapping is intentionally surface-oriented: secure messaging protects the
external APDU-facing boundary, Security Domains express administrative and
cryptographic authority, the kernel privilege model mediates privileged
effects, and registry consistency plus runtime isolation limit the impact of a
hostile component once code has been deployed inside the platform.

Depending on its mode, the secure channel protocol provides peer authentication,
message integrity, confidentiality, and stateful replay protection for APDU
exchanges. These properties can prevent an unauthenticated endpoint from
modifying or impersonating a protected session, but do not prevent a transparent
real-time relay without an additional proximity or channel-binding mechanism
\cite{roland2012relay}.

The privilege model is intended to restrict administrative operations to the
authority assigned to the active administrative context, addressing privilege
escalation threats observed in malicious application scenarios
\cite{mostowski2008malicious}.

Security Domains express the delegation of administrative and cryptographic
authority over application subtrees. They do not themselves enforce memory or
execution isolation: those boundaries are provided by the kernel and, where
available, the processor protection mechanisms.

Registry consistency mechanisms are designed to validate state transitions and
preserve a recoverable valid state when committing persistent changes, thereby
mitigating attacks and failures targeting internal metadata
\cite{munoz2023tee}.

Runtime isolation is intended to prevent applications from accessing or
interfering with each other's execution or data, a critical defense against
malicious applications \cite{busch2020trustedcore}.

%--------------------------------------------------
\section{Out-of-Scope Considerations}

This section explicitly defines what is not covered by the threat model.

\begin{itemize}
\item Physical attacks
\item Side-channel attacks
\item Hardware compromise
\item Correctness of an application's own behaviour
\item Security claims for QEMU-oriented builds
\end{itemize}

The out-of-scope boundary is particularly important for the hardware-origin
attack surface introduced in Section~2.1: hardware compromise and lower-level
silicon-origin attacks are explicitly excluded from the present chapter because
they could invalidate the confidentiality and integrity properties sought by
service providers.

Physical attacks, such as fault injection or probing, are not considered in
this model, although they are well studied in the literature
\cite{boneh1997faults}.

Side-channel attacks, including timing or power analysis, are also excluded,
despite their practical impact \cite{kocher1996timing, kocher1999dpa}.

Hardware compromise falls outside the scope of this architecture-level model.

Application-level correctness is not an assumption made by this threat model:
an ordinary Rustlet may be defective or malicious. Protecting that Rustlet's
own assets from its own implementation bugs is out of scope, while containing
their effects on other applications and on kernel-owned state is a platform
security objective. The same distinction applies to a pre-deployed ordinary
Rustlet; pre-issuance placement does not turn its developer into a trusted
platform authority.

No deployment-security claim is made for QEMU-oriented builds. They are
retained as development and validation tools, not as secure deployment targets.
This is particularly important because the project
uses semihosting during bring-up and debugging, and semihosting explicitly
introduces a host-assisted escape hatch that breaks the isolation assumptions
expected from a secure element style environment \cite{qemu-security,
qemu-semihosting}. The threat model therefore applies to intended hardware
deployments, not to the convenience execution environments used during
prototype development.

%--------------------------------------------------
